\documentclass[11pt]{article}

\usepackage[margin=1in]{geometry}
\usepackage{amsmath,amssymb,amsthm,mathtools}
\usepackage{array,booktabs,longtable}
\usepackage[T1]{fontenc}
\usepackage{lmodern}
\usepackage{microtype}
\usepackage{xcolor}
\usepackage{xurl}
\usepackage[colorlinks=true,linkcolor=blue!55!black,
            urlcolor=blue!55!black]{hyperref}

\newtheorem{theorem}{Theorem}[section]
\newtheorem{proposition}[theorem]{Proposition}
\newtheorem{lemma}[theorem]{Lemma}
\newtheorem{corollary}[theorem]{Corollary}
\theoremstyle{remark}
\newtheorem{remark}[theorem]{Remark}

\newcommand{\Z}{\mathbf Z}
\newcommand{\F}{\mathbf F}
\newcommand{\Gm}{\mathbf G_m}
\newcommand{\Ga}{\mathbf G_a}
\newcommand{\Spec}{\operatorname{Spec}}
\newcommand{\Soc}{\operatorname{Soc}}
\newcommand{\gr}{\operatorname{gr}}
\newcommand{\length}{\operatorname{length}}
\newcommand{\eps}{\varepsilon}
\newcommand{\m}{\mathfrak m}
\newcommand{\n}{\mathfrak n}
\newcommand{\ot}{\otimes}
\newcommand{\id}{\operatorname{id}}

\title{Audit and Conceptual Simplification of a Length-Nine\\
Rank-Four Group Scheme Not Killed by Four}
\author{Independent algebraic and computational verification}
\date{11 July 2026}

\begin{document}
\maketitle

\begin{abstract}
We audit the claimed rank-four example over a length-nine Artin ring and find
that it is correct.  The base is the two-generator Artin complete intersection
\[
 R=(\Z/4)[x,y]/(x^3,y^3,xy^2-2),
\]
and the coordinate algebra is the rank-four relative complete intersection
\[
 A=R[U,V]/(U^2-xyU-x^2V,\;V^2-y^2V).
\]
Writing $W=UV$ and
\[
 \lambda=1+yU+xV+xyW=(1+yU)(1+xV),
\]
the coproduct is compressed into the two skew-primitive formulas
\[
 \Delta(U)=U\ot1+\lambda\ot U,
 \qquad
 \Delta(V)=V\ot\lambda+1\ot V.
\]
We give hand proofs of the base length, freeness, every Hopf axiom, and the
power calculation.  The shortest explanation is a norm calculation.  If
$N_n(T)=1+T+\cdots+T^{n-1}$, then
\[
 [n]^\#(U)=N_n(\lambda)U,
 \qquad [n]^\#(V)=VN_n(\lambda).
\]
Here $\lambda=1+\theta$, $\theta^2=0$, and $2\theta=2xV$.  Hence
\[
 N_4(\lambda)=2\theta=2xV,
 \qquad [4]^\#(U)=2xUV\ne0,
 \qquad [8]=e.
\]
Thus the example is a characteristic-four ramified norm phenomenon inside an
ambient semidirect product: a binomial carry turns the fourth power into the
top socle class.
The presentation is already generator-minimal in the natural sense, but this
norm description makes it substantially simpler conceptually.  Length nine
is minimal for this group-like bridge family and for quotients of this base;
global nonexistence at length eight is not proved.
\end{abstract}

\begin{center}
\fbox{\begin{minipage}{0.92\textwidth}
\small
\textbf{Disclosure: AI-generated work.}\ This document in its entirety---the
construction, the proofs, the computations, and the exposition---was generated
by artificial intelligence. The work was carried out by the AI models Codex
(OpenAI) and Claude (Anthropic).
\end{minipage}}
\end{center}
\medskip

\tableofcontents

\section{Verdict and exact statement}

The claim is correct.  It is useful to write the base as
\begin{equation}\label{eq:R}
 R=(\Z/4)[x,y]/(x^3,y^3,xy^2-2).
\end{equation}
The original presentation over $\Z$ is identical: the three relations
$x^3=y^3=0$ and $xy^2=2$ already imply
\[
 4=(xy^2)^2=x^2y^4=0.
\]
Define
\begin{equation}\label{eq:A}
 A=R[U,V]/(U^2-xyU-x^2V,\;V^2-y^2V),
 \qquad W=UV,
\end{equation}
and put
\begin{equation}\label{eq:theta-lambda}
 \theta=yU+xV+xyW,
 \qquad
 \lambda=1+\theta=(1+yU)(1+xV).
\end{equation}

\begin{theorem}\label{thm:main}
The ring $R$ is local Artin Gorenstein with residue field $\F_2$,
characteristic $4$, cardinality $2^9$, length $9$, Hilbert function
$(1,2,3,2,1)$, and socle $\F_2(2x)$.

The $R$-algebra $A$ is free with basis $1,U,V,W$.  There is a Hopf algebra
structure with $\eps(U)=\eps(V)=0$ and
\begin{equation}\label{eq:Delta}
 \Delta(U)=U\ot1+\lambda\ot U,
 \qquad
 \Delta(V)=V\ot\lambda+1\ot V.
\end{equation}
For $G=\Spec A$, the fourth-power word is not the unit word, whereas the
eighth-power word is:
\begin{equation}\label{eq:powers-main}
 [4]^\#(U)=2xW\ne0,
 \qquad [4]^\#(V)=[4]^\#(W)=0,
 \qquad [8]=e.
\end{equation}
Moreover $G$ is noncommutative.
\end{theorem}

For a noncommutative group scheme, $[n]$ denotes the power word
$g\mapsto g^n$; it need not be a group homomorphism.  The phrase ``killed by
$n$'' means that this word factors through the unit section.  Contravariantly,
the constant unit word has comorphism
\[
 A\xrightarrow{\eps}R\xrightarrow{\iota}A.
\]
Thus the nonzero coordinate in \eqref{eq:powers-main} proves exactly the
claimed failure, with no appeal to the existence of an ordinary $R$-valued
point witnessing it.

\subsection*{A first explanation using only basic Hopf algebra}
\addcontentsline{toc}{subsection}{A first explanation using only basic Hopf algebra}

Here is a way to discover the construction backwards from the desired
conclusion.  It uses only the elementary dictionary between an affine group
scheme and its coordinate Hopf algebra.

Suppose $G=\Spec A$.  The group multiplication corresponds to the coproduct
\[
 \Delta:A\longrightarrow A\ot_R A.
\]
The squaring map $g\mapsto g^2$ is obtained by applying the group
multiplication to the diagonal.  On coordinate rings this means
\begin{equation}\label{eq:first-reading-square}
 [2]^\#=m_A\circ\Delta,
\end{equation}
where $m_A:A\ot_R A\to A$ is ordinary multiplication.

Now recall two standard definitions.  An element $\lambda\in A$ is
\emph{group-like} if
\[
 \Delta(\lambda)=\lambda\ot\lambda.
\]
An element $U\in A$ is left $\lambda$-\emph{skew primitive} if
\[
 \Delta(U)=U\ot1+\lambda\ot U.
\]
Applying \eqref{eq:first-reading-square} immediately gives
\[
 [2]^\#(U)=(1+\lambda)U.
\]
Repeating the same calculation gives the geometric-series formula
\begin{equation}\label{eq:first-reading-norm}
 [n]^\#(U)=N_n(\lambda)U,
 \qquad
 N_n(T)=1+T+\cdots+T^{n-1}.
\end{equation}
This formula is the central idea.  To make $[4]^\#(U)$ nonzero, we need to
make the product $N_4(\lambda)U$ nonzero.

The characteristic-four trick is to take
\[
 \lambda=1+\theta,
 \qquad \theta^2=0.
\]
Then $\lambda^i=1+i\theta$, so
\begin{equation}\label{eq:first-reading-carry}
 N_4(\lambda)
 =1+(1+\theta)+(1+2\theta)+(1+3\theta)
 =4+6\theta=2\theta.
\end{equation}
Thus $\lambda^4=1$, but its fourth geometric sum need not vanish.  The whole
problem has now been reduced to the following very concrete task:
\begin{equation}\label{eq:first-reading-task}
 \text{construct a rank-four Hopf algebra with}
 \quad \theta^2=0
 \quad\text{and}\quad 2\theta U\ne0.
\end{equation}

A convenient way to keep $2\theta U$ alive is to introduce a second
coordinate $V$.  Take
\[
 R=(\Z/4)[x,y]/(x^3,y^3,xy^2-2)
\]
and
\[
 A=R[U,V]/(U^2-xyU-x^2V,\;V^2-y^2V).
\]
Because both relations are monic quadratics, reduction of powers of $U$ and
$V$ leaves exactly the four-element $R$-basis
\begin{equation}\label{eq:first-reading-basis}
 1,\qquad U,\qquad V,\qquad UV.
\end{equation}
This is why the resulting group scheme has rank four.

Define
\begin{equation}\label{eq:first-reading-lambda}
 \theta=yU+xV+xyUV,
 \qquad
 \lambda=1+\theta=(1+yU)(1+xV).
\end{equation}
The relations give
\begin{equation}\label{eq:first-reading-identities}
 \theta^2=0,
 \qquad
 2\theta=2xV.
\end{equation}
The coproduct is chosen to make $\lambda$ group-like and to make the two
coordinates skew primitive in opposite directions:
\begin{equation}\label{eq:first-reading-coproduct}
 \Delta(U)=U\ot1+\lambda\ot U,
 \qquad
 \Delta(V)=V\ot\lambda+1\ot V.
\end{equation}
The same geometric-series argument applies to the right skew-primitive
element $V$; the order of the factors causes no problem because the
coordinate algebra $A$ is commutative.
The later Hopf-algebra audit verifies directly that these formulas respect
the two quadratic relations and that
$\Delta(\lambda)=\lambda\ot\lambda$.  Once those checks are made, the power
calculation is automatic from \eqref{eq:first-reading-norm}.

Indeed, \eqref{eq:first-reading-carry} and
\eqref{eq:first-reading-identities} give
\begin{equation}\label{eq:first-reading-fourth}
 [4]^\#(U)=2\theta U=2xUV,
 \qquad
 [4]^\#(V)=V(2\theta)=2xV^2=4V=0.
\end{equation}
The first answer is nonzero.  The relation $xy^2=2$ gives
\[
 2x=x^2y^2\ne0
\]
in $R$, and $UV$ is one of the free basis elements in
\eqref{eq:first-reading-basis}.  Therefore $2xUV\ne0$ in $A$.

Finally, $\theta^2=0$ and $4=0$ imply $\lambda^4=1$, and hence
\[
 N_8(\lambda)
 =N_4(\lambda)(1+\lambda^4)
 =2N_4(\lambda)=0.
\]
The norm formula therefore sends both $U$ and $V$ to zero at the eighth
power.  Since $U,V$ generate the augmented $R$-algebra $A$, $[8]$ is the
unit word.  The complete construction can therefore be
remembered as the single chain
\begin{equation}\label{eq:first-reading-box}
 \boxed{
 \begin{gathered}
 \Delta(U)=U\ot1+\lambda\ot U
 \quad\Longrightarrow\quad
 [4]^\#(U)=N_4(\lambda)U,\\
 \lambda=1+\theta,\quad\theta^2=0
 \quad\Longrightarrow\quad N_4(\lambda)=2\theta,\\
 2\theta=2xV
 \quad\Longrightarrow\quad [4]^\#(U)=2xUV\ne0.
 \end{gathered}}
\end{equation}

In words: a group-like element infinitesimally close to $1$ leaves a
nonzero geometric-series carry in characteristic four, and the second
coordinate $V$ stores that carry long enough for multiplication by $U$ to
turn it into the surviving top class $2xUV$.  The remaining sections prove
all the ring and Hopf-algebra identities used in this short explanation.  A
later section also interprets the same calculation inside an ordinary
semidirect product.

\subsection*{A deliberately unoptimized five-coefficient construction}
\addcontentsline{toc}{subsection}{A deliberately unoptimized five-coefficient construction}

The preceding explanation still presents the optimized coefficients
$xy,x^2,y^2,y,x$.  If the goal is to see why a construction of this shape
exists, rather than to minimize the base, it is clearer not to solve for
those coefficients.  Treat the five coefficients as separate generators.

The elementary building block is worth isolating first.  Suppose
$a,b$ belong to a ring $D$ and $ab=-2$.  On
\[
 C=D[T]/(T^2-aT),
 \qquad g=1+bT,
\]
try the skew-primitive formula
\[
 \Delta(T)=T\ot1+g\ot T.
\]
If $T_*=T_1+g_1T_2$, then one line of calculation gives
\begin{equation}\label{eq:rank-two-block}
 T_*^2-aT_*
 =g_1\bigl(2T_1+a(g_1-1)\bigr)T_2
 =g_1(2+ab)T_1T_2=0.
\end{equation}
Moreover $\Delta(g)=g\ot g$ and
$g^2=1+b(2+ab)T=1$.  Thus the single scalar equation $ab=-2$
produces a rank-two Hopf algebra: take $\eps(T)=0$ and $S(T)=T$;
group-likeness of $g$ makes coassociativity formal, and
$T+gS(T)=0=S(T)+S(g)T$.  In other words, the term $abT_1T_2$
cancels the cross-term $2T_1T_2$.  The whole construction can be checked by
someone who has only just learned the definition of a Hopf algebra.

The rank-four construction begins with two such cancellation patterns,
$(\alpha,r)$ and $(\gamma,s)$.  It then inserts one off-diagonal
``bridge'' term $\beta V$ into the equation for $U^2$.  Three extra scalar
equations make that bridge compatible with the two rank-two blocks.  The
unoptimized recipe is therefore simply
\begin{equation}\label{eq:large-five-relations}
 \underbrace{\alpha r=-2,\quad \gamma s=-2}_{
   \text{two rank-two patterns}},
 \qquad
 \underbrace{\gamma r=0,\quad\beta s=0,\quad
   \beta r=-\alpha s}_{\text{one compatible bridge}}.
\end{equation}
Rather than solve these equations economically, we use them as the
definition of the base ring.

\begin{equation}\label{eq:large-base}
 B_{\mathrm{br}}=\Z[\alpha,\beta,\gamma,r,s]
   /\bigl(\alpha r+2,\;\gamma s+2,\;
           \gamma r,\;\beta s,\;\beta r+\alpha s\bigr).
\end{equation}
We make no attempt to force $B_{\mathrm{br}}$ to be Artinian: none is
needed.  Indeed, setting $\alpha=\beta=\gamma=r=0$ gives a surjection
$B_{\mathrm{br}}\twoheadrightarrow\F_2[s]$.  The group scheme will nevertheless be
finite because its coordinate algebra is free of rank four over
$B_{\mathrm{br}}$.

Over $B_{\mathrm{br}}$, put
\begin{equation}\label{eq:large-A}
 A_{\mathrm{br}}=B_{\mathrm{br}}[U,V]
 /(U^2-\alpha U-\beta V,\;V^2-\gamma V),
 \qquad W=UV,
\end{equation}
and
\begin{equation}\label{eq:large-lambda}
 \lambda=(1+rU)(1+sV),
 \qquad
 \theta=\lambda-1.
\end{equation}
Define
\begin{equation}\label{eq:large-Delta}
 \Delta(U)=U\ot1+\lambda\ot U,
 \qquad
 \Delta(V)=V\ot\lambda+1\ot V,
 \qquad
 \eps(U)=\eps(V)=0.
\end{equation}

\begin{proposition}[The unoptimized bridge]\label{prop:large-bridge}
The formulas \eqref{eq:large-A}--\eqref{eq:large-Delta} make
$A_{\mathrm{br}}$ a finite free Hopf algebra of rank four over
$B_{\mathrm{br}}$.  Its power words satisfy
\begin{equation}\label{eq:large-powers}
 [4]^\#(U)=2sW\ne0,
 \qquad [4]^\#(V)=[4]^\#(W)=0,
 \qquad [8]=e.
\end{equation}
\end{proposition}

\begin{proof}
The relations in \eqref{eq:large-five-relations} imply
\begin{equation}\label{eq:large-consequences}
 4=0,\qquad
 2\alpha=2\beta=2\gamma=2r=0,\qquad
 \beta r^2=2s,\qquad 2s^2=0.
\end{equation}
For example,
\[
 4=(\alpha r)(\gamma s)=\alpha s(\gamma r)=0,
\]
while
\[
 2\gamma=-(\alpha r)\gamma=0,
 \quad 2r=-(\gamma s)r=0,
 \quad 2\beta=-(\gamma s)\beta=0,
\]
and
\[
 2\alpha=-(\gamma s)\alpha
          =-\gamma(\alpha s)=\gamma\beta r=0.
\]
Multiplying $\beta r+\alpha s=0$ by $r$ gives
$\beta r^2=2s$, and then multiplication by $s$ gives $2s^2=0$.
Since the defining relations of $A_{\mathrm{br}}$ are successively
monic in $V$ and $U$, respectively,
\[
 A_{\mathrm{br}}=B_{\mathrm{br}}\{1,U,V,W\}
\]
is free of rank four.

The identity $\beta r^2=2s$ is also the reason the bridge is essential.
If one set $\beta=0$, then it would force $2s=0$, and the fourth-power
class computed below would disappear.  The bridge is precisely what keeps
the characteristic-four carry alive.

We next explain why the Hopf structure works without listing a large
coefficient table.  Put
\[
 L=1+rU,\qquad M=1+sV,\qquad\lambda=LM.
\]
Using \eqref{eq:large-five-relations} and the two quadratic relations gives
\[
 M^2=1,\qquad L^2=1+2sV,\qquad
 \lambda^2=1+2sV.
\]
Also $2\theta=2sV$, so comparison with
$(1+\theta)^2$ gives
\begin{equation}\label{eq:large-theta}
 \theta^2=0,\qquad
 2\theta=2sV,\qquad
 \lambda^{-1}=1-\theta.
\end{equation}

Let $U_i,V_i,\lambda_i$ denote the coordinates in the two tensor factors and
set
\[
 U_*=U_1+\lambda_1U_2,
 \qquad
 V_*=V_1\lambda_2+V_2.
\]
The five relations in \eqref{eq:large-five-relations} give the closure
identities
\begin{equation}\label{eq:large-closure}
 U_*^2=\alpha U_*+\beta V_*,
 \qquad
 V_*^2=\gamma V_*,
 \qquad
 (1+rU_*)(1+sV_*)=\lambda_1\lambda_2.
\end{equation}
Here is a short view of the cancellations.  One has
\[
 \gamma(\lambda-1)=-2V,\qquad
 \beta(1-\lambda)=\alpha sU,\qquad
 \beta(\lambda^2-1)=0,
\]
and
\[
 2\lambda U+\alpha(\lambda^2-\lambda)+\alpha sV=0.
\]
Indeed, they give the transparent defect calculations
\[
 V_*^2-\gamma V_*
 =V_1\lambda_2\bigl(\gamma(\lambda_2-1)+2V_2\bigr)=0
\]
and
\begin{align*}
 U_*^2-\alpha U_*-\beta V_*
  ={}&\beta V_1(1-\lambda_2)+2\lambda_1U_1U_2
      +\alpha(\lambda_1^2-\lambda_1)U_2
      +\beta(\lambda_1^2-1)V_2\\
  ={}&\bigl(\alpha sV_1+2\lambda_1U_1
      +\alpha(\lambda_1^2-\lambda_1)\bigr)U_2=0.
\end{align*}
For the group-like identity, put
$E=rsV_1U_2$.  Direct factorization gives
\begin{equation}\label{eq:large-group-like-factor}
 (1+rU_*)(1+sV_*)-\lambda_1\lambda_2
  =(L_1M_2)E\bigl(2+rU_2+sV_1+E\bigr).
\end{equation}
Each of the four terms on the right vanishes by the two quadrics and the
consequences of \eqref{eq:large-five-relations}: $2E=0$,
$E(rU_2)=0$, $E(sV_1)=0$, and $E^2=0$.
Thus \eqref{eq:large-closure} holds.  It says precisely that $\Delta$
respects the two quadrics and that $\lambda$ is group-like.  Counitality is
immediate, and group-likeness makes coassociativity of the two skew-primitive
formulas formal.

An antipode is
\begin{equation}\label{eq:large-antipode}
 S(U)=U-sW+\alpha sV,
 \qquad
 S(V)=V+rW.
\end{equation}
Indeed these are $-\lambda^{-1}U$ and $-V\lambda^{-1}$, respectively.
For a direct check, use
\[
 UW=\alpha W+\beta\gamma V,\qquad
 VW=\gamma W,\qquad
 W^2=\alpha\gamma W+\beta\gamma^2V.
\]
Substitution gives
\[
 S(U)^2=\alpha S(U)+\beta S(V),\qquad
 S(V)^2=\gamma S(V),\qquad S(\lambda)=\lambda^{-1}.
\]
Thus $S$ defines an algebra map, and the two convolution identities follow
from \eqref{eq:large-Delta} and the formulas
$S(U)=-\lambda^{-1}U$, $S(V)=-V\lambda^{-1}$.

It remains to compute the power word.  The norm formula
\eqref{eq:first-reading-norm} and \eqref{eq:large-theta} give
\[
 N_4(\lambda)=4+6\theta=2\theta=2sV.
\]
Therefore
\[
 [4]^\#(U)=2sVU=2sW,
 \qquad
 [4]^\#(V)=2sV^2=2s\gamma V=-4V=0.
\]
Since the power word is still a morphism of schemes, its comorphism is an
algebra map, and hence
\[
 [4]^\#(W)=[4]^\#(U)[4]^\#(V)=0.
\]
To see directly that the first element is nonzero, use the specialization
\begin{equation}\label{eq:large-specialization-in-proof}
 B_{\mathrm{br}}\longrightarrow R,\qquad
 (r,s,\alpha,\beta,\gamma)\longmapsto(y,x,xy,x^2,y^2).
\end{equation}
The five defining equations of $B_{\mathrm{br}}$ vanish in $R$, and $2s$
maps to the nonzero element $2x$ (see \eqref{eq:carry}).  Explicitly, the
first two equations map to $4=0$,
the next two to $y^3=x^3=0$, and the last to
$2x^2y=x^3y^3=0$.  Thus $2s\ne0$ in $B_{\mathrm{br}}$, and freeness of
$A_{\mathrm{br}}$ gives
$2sW\ne0$.  Finally
$\lambda^4=1$ and
\[
 N_8(\lambda)=N_4(\lambda)(1+\lambda^4)=2N_4(\lambda)=0,
\]
so $[8]=e$.
\end{proof}

\subsubsection*{A human-sized untruncated base}

The five-coefficient ring explains where the equations come from.  For a
short concrete presentation which is still untruncated, solve only two of
the coefficients,
\[
 \alpha=rs,\qquad \beta=s^2,
\]
and leave $\gamma=g$ as a separate generator.  This gives
\begin{equation}\label{eq:three-coefficient-base}
 D=\Z[r,s,g]/(r^2s+2,\;s^3,\;gs+2,\;gr).
\end{equation}
The four relations are easy to remember: $(rs,r)$ and $(g,s)$ are the two
rank-two cancellation patterns, while $s^3=gr=0$ kill the two cross-terms.
The remaining bridge balance is automatic.  Indeed, $(gs+2)r=0$ and
$gr=0$ give $2r=0$, whence
\[
 s^2r+rs^2=2rs^2=0.
\]
Consequently there is a surjection
\begin{equation}\label{eq:bridge-to-D}
 B_{\mathrm{br}}\twoheadrightarrow D,
 \qquad
 (\alpha,\beta,\gamma,r,s)\longmapsto(rs,s^2,g,r,s).
\end{equation}

Thus Proposition~\ref{prop:large-bridge} immediately gives the following
particularly short example:
\begin{equation}\label{eq:three-coefficient-A}
 A_D=D[U,V]/(U^2-rsU-s^2V,\;V^2-gV),
 \qquad \lambda=(1+rU)(1+sV),
\end{equation}
with
\begin{equation}\label{eq:three-coefficient-Delta}
 \Delta(U)=U\ot1+\lambda\ot U,
 \qquad
 \Delta(V)=V\ot\lambda+1\ot V.
\end{equation}
It is free over $D$ on $1,U,V,UV$, and
\begin{equation}\label{eq:three-coefficient-powers}
 [4]^\#(U)=2sUV\ne0,
 \qquad [4]^\#(V)=[4]^\#(UV)=0,
 \qquad [8]=e.
\end{equation}
Here nonvanishing follows from the specialization
\[
 D\twoheadrightarrow R,
 \qquad (r,s,g)\longmapsto(y,x,y^2),
\]
which sends $2s$ to $2x\ne0$.  The base $D$ is genuinely untruncated and
non-Artinian: setting $2=s=g=0$ gives a surjection
$D\twoheadrightarrow\F_2[r]$.

The relationship among the three presentations is now completely visible:
\[
 B_{\mathrm{br}}\twoheadrightarrow D\twoheadrightarrow R,
 \qquad
 D/(g-r^2)\simeq
 \Z[r,s]/(r^3,s^3,r^2s+2)\simeq R.
\]
The five-coefficient ring is best for deriving the compatibility equations;
$D$ is arguably the easiest concrete example; and $R$ is the optimized
length-nine specialization.

No minimality claim is being made for either enlarged ansatz.  The later
statement is narrower: after imposing
$\alpha=rs$, $\beta=s^2$, and $\gamma=r^2$, the Hopf identities force a map
from the universal coefficient ring $R$, and survival of the fourth-power
class forces that map to be injective.  Thus the exact two-coefficient
family contains a copy of the length-nine ring, whereas the constructions
above deliberately leave some or all of $\alpha,\beta,\gamma$ unoptimized.

\section{The length-nine base without guesswork}

\subsection{A Smith-normal-form proof}

There is a particularly clean proof that the integer presentation really has
length nine.  Let
\[
 B=\Z[x,y]/(x^3,y^3).
\]
It is free of rank nine over $\Z$, on the monomials $x^iy^j$ with
$0\le i,j\le2$.  Let $N:B\to B$ be multiplication by $xy^2$.  Then
\[
 R_{\mathrm{add}}=\operatorname{coker}(N-2\id).
\]
On the nine monomials, $N$ has exactly two nonzero arrows:
\[
 1\longmapsto xy^2,
 \qquad x\longmapsto x^2y^2.
\]
It is zero on the other five complementary basis elements and $N^2=0$.
Each of the two $2\times2$ blocks of $N-2\id$ has Smith normal form
$\operatorname{diag}(1,4)$, while each of the five remaining blocks is
$(-2)$.  Consequently
\begin{equation}\label{eq:additive-R}
 R_{\mathrm{add}}\simeq(\Z/4)^2\oplus(\Z/2)^5.
\end{equation}
More explicitly, every element has a unique normal form
\begin{equation}\label{eq:normal-form-R}
 a+bx+\alpha y+\beta x^2+\gamma xy+\delta y^2+\eta x^2y,
\end{equation}
where $a,b\in\Z/4$ and
$\alpha,\beta,\gamma,\delta,\eta\in\F_2$.

This calculation proves at once that $|R|=2^9$, that $2\ne0$, and that
\begin{equation}\label{eq:carry}
 xy^2=2,
 \qquad x^2y^2=2x\ne0.
\end{equation}
It also proves that multiplication by $xy^2-2$ on $B$ is injective: its
determinant is $(-2)^9\ne0$.  Therefore
\[
 x^3,\quad y^3,\quad xy^2-2
\]
is a regular sequence in $\Z[x,y]$.  Thus $R$ is an Artin complete
intersection, which gives an independent conceptual proof of its
Gorenstein property.

\subsection{Locality, filtration, and socle}

The relations imply
\begin{equation}\label{eq:torsion}
 2y=2x^2=2xy=2y^2=2x^2y=0.
\end{equation}
Both $x$ and $y$ are nilpotent, so every prime ideal contains them; it then
also contains $2=xy^2$.  Hence
\[
 \m=(x,y)
\]
is the unique maximal ideal and $R/\m=\F_2$.  This nilpotence observation is
needed for uniqueness; merely noting that $R/(x,y)=\F_2$ would prove only
that $(x,y)$ is maximal.

The associated graded ring is especially simple:
\begin{equation}\label{eq:grR}
 \gr_{\m}R\simeq\F_2[X,Y]/(X^3,Y^3).
\end{equation}
Indeed, $x,y$ generate $\m$, so there is a surjection from the right-hand
side to $\gr_{\m}R$.  Both sides have $\F_2$-dimension nine by
\eqref{eq:additive-R}, so it is an isomorphism.  In concrete terms,
\begin{align}
 \m^2&=\F_2\{x^2,xy,y^2,x^2y,2,2x\},\nonumber\\
 \m^3&=\F_2\{x^2y,2,2x\},\nonumber\\
 \m^4&=\F_2\{2x\},
 &\m^5&=0.\label{eq:m-powers}
\end{align}
Thus the Hilbert function is $(1,2,3,2,1)$ and the length is nine.  Direct
multiplication of the normal form \eqref{eq:normal-form-R} by $x$ and $y$
gives
\begin{equation}\label{eq:socle}
 \Soc(R)=\F_2(2x).
\end{equation}
In the graded picture, this is simply the top monomial
$X^2Y^2$, lifted by $x^2y^2=2x$.

\section{The rank-four algebra and the bridge element}

First form the monic quadratic extension
\[
 C=R[V]/(V^2-y^2V),
\]
which is free on $1,V$.  Then
\[
 A=C[U]/(U^2-xyU-x^2V)
\]
is free on $1,U$ over $C$.  This proves, without a Groebner-basis
assumption, that
\begin{equation}\label{eq:A-basis}
 A=R\{1,U,V,W\},\qquad W=UV,
\end{equation}
is finite locally free of rank four.

The complete multiplication needed below is
\begin{equation}\label{eq:A-multiplication}
\begin{aligned}
 U^2&=xyU+x^2V,& V^2&=y^2V,& UV&=W,\\
 UW&=xyW+2xV,& VW&=y^2W,& W^2&=0.
\end{aligned}
\end{equation}
For example,
\[
 UW=U^2V=xyW+x^2y^2V=xyW+2xV,
\]
and the remaining identities follow similarly.

\begin{lemma}\label{lem:theta}
For $\theta$ and $\lambda$ in \eqref{eq:theta-lambda},
\begin{equation}\label{eq:theta-identities}
 \theta^2=0,
 \qquad \lambda^{-1}=1-\theta,
 \qquad 2\theta=2xV,
 \qquad \lambda^2=1+2xV,
 \qquad \lambda^4=1.
\end{equation}
\end{lemma}

\begin{proof}
The two nonzero-looking square terms in $\theta^2$ are
\[
 y^2U^2=2xV,
 \qquad x^2V^2=2xV.
\]
Their sum is $4xV=0$.  All mixed terms vanish by
\eqref{eq:torsion}, $x^3=y^3=0$, or $W^2=0$.  Thus $\theta^2=0$.
The inverse formula follows.  Also $2y=2xy=0$, so
$2\theta=2xV$.  Finally
\[
 (1+\theta)^2=1+2\theta=1+2xV
\]
and its square is one.
\end{proof}

Modulo $\m$, one has
\[
 A/\m A\simeq\F_2[U,V]/(U^2,V^2).
\]
Thus two coordinate generators are unavoidable: the cotangent space at the
unit of the special fiber has dimension two.

\section{A hand audit of the Hopf structure}

The only delicate questions are whether \eqref{eq:Delta} respects the two
quadrics and whether $\lambda$ is group-like.  We verify both directly.

In $A\ot_R A$, use lower-case abbreviations
\[
 u=U\ot1,\quad v=V\ot1,\quad
 u'=1\ot U,\quad v'=1\ot V,
\]
and put $l=\lambda\ot1=1+\theta_1$ and
$l'=1\ot\lambda=1+\theta_2$.  The proposed product coordinates are
\begin{equation}\label{eq:product-coordinates}
 \bar u=u+lu',
 \qquad \bar v=vl'+v'.
\end{equation}

\subsection{Preservation of the two quadrics}

Using the defining relations in each tensor factor gives
\begin{align}
 \bar u^2-xy\bar u-x^2\bar v
  ={}&x^2v(1-l')+2luu'+xy(l^2-l)u'
      +x^2(l^2-1)v'\nonumber\\
  ={}&-x^2v\theta_2+2(1+\theta_1)uu'
      +xy\theta_1u'+2x^2\theta_1v'.\label{eq:U-defect}
\end{align}
The following four reductions are immediate from
\eqref{eq:A-multiplication} and the base relations:
\[
 x^2v\theta_2=x^2yv u',
 \qquad 2\theta_1u=2xuv,
\]
\[
 xy\theta_1=2u+x^2yv+2xuv,
 \qquad 2x^2\theta_1=0.
\]
Substitution in \eqref{eq:U-defect} gives
\[
 -x^2yvu'+(2uu'+2xuvu')
 +(2uu'+x^2yvu'+2xuvu')=0.
\]

Similarly,
\begin{align}
 \bar v^2-y^2\bar v
 &=y^2v(l'^2-l')+2vv'l'\nonumber\\
 &=y^2v\theta_2+2vv'(1+\theta_2).
 \label{eq:V-defect}
\end{align}
Here $y^2v\theta_2=2vv'$ and $2vv'\theta_2=0$, so
\eqref{eq:V-defect} is $4vv'=0$.  Hence the proposed coproduct descends to
the quotient algebra $A$.

\subsection{The group-like identity}

Write
\[
 a_1=1+yu,\quad b_1=1+xv,
 \qquad a_2=1+yu',\quad b_2=1+xv'.
\]
Then $l=a_1b_1$ and $l'=a_2b_2$.  From
\eqref{eq:product-coordinates},
\[
 1+y\bar u=a_1(1+yb_1u'),
 \qquad
 1+x\bar v=b_2(1+xva_2).
\]
Set $p=xv$ and $q=yu'$.  The difference between
\[
 (1+q(1+p))(1+p(1+q))
 \quad\text{and}\quad (1+p)(1+q)
\]
is
\[
 pq(2+p+q+pq).
\]
But $2pq=p^2q=pq^2=0$, by \eqref{eq:torsion} and the two quadrics.
It follows that
\begin{equation}\label{eq:group-like}
 \Delta(\lambda)=\lambda\ot\lambda.
\end{equation}

For reference, in rows and columns indexed by $1,U,V,W$, the same identity
is the compact matrix certificate
\[
 M(\Delta\lambda)=
 \begin{pmatrix}
 1&y&x&xy\\
 y&y^2&xy&2\\
 x&xy&x^2&x^2y\\
 xy&2&x^2y&2x
 \end{pmatrix}
 =
 \begin{pmatrix}1\\y\\x\\xy\end{pmatrix}
 \begin{pmatrix}1&y&x&xy\end{pmatrix}.
\]

\subsection{Counit, coassociativity, and antipode}

The augmentation $\eps(U)=\eps(V)=0$ is plainly a counit.  Once
\eqref{eq:group-like} is known, coassociativity is formal.  For example,
\begin{align*}
 (\Delta\ot\id)\Delta(U)
 &=U\ot1\ot1+\lambda\ot U\ot1
   +\lambda\ot\lambda\ot U\\
 &=(\id\ot\Delta)\Delta(U),
\end{align*}
and the same calculation with the skewness reversed applies to $V$.

The antipode can be written either conceptually or in the basis
\eqref{eq:A-basis}:
\begin{align}
 S(U)&=-\lambda^{-1}U=U+x^2yV-xW,\nonumber\\
 S(V)&=-V\lambda^{-1}=V+yW,
 &S(W)&=W.\label{eq:antipode}
\end{align}
Substitution in \eqref{eq:A-multiplication} shows that these images respect
both quadrics and that $S(\lambda)=\lambda^{-1}$.  The two convolution
identities on $U$ and $V$ now read
\[
 S(U)+S(\lambda)U=0,
 \qquad U+\lambda S(U)=0,
\]
\[
 S(V)\lambda+V=0,
 \qquad VS(\lambda)+S(V)=0.
\]
They hold on all of $A$ because the displayed maps are algebra maps into the
commutative algebra $A$.  This completes a hand proof of every Hopf axiom.

Finally, $G$ is noncommutative.  The coefficient of $V\ot U$ in
$\Delta(U)$ is $x$, whereas the coefficient of $U\ot V$ is zero.  Since
$x\ne0$ and $A\ot_R A$ is free, $\Delta$ is not cocommutative.

\section{The conceptual power calculation}

The following elementary lemma is the main simplification of the example.

\begin{lemma}[Skew-primitive norm formula]\label{lem:norm}
Let $B$ be a commutative Hopf algebra, let $l\in B$ be group-like, and
suppose
\[
 \Delta(z)=z\ot1+l\ot z
 \quad\text{or}\quad
 \Delta(z)=z\ot l+1\ot z.
\]
For $N_n(T)=1+T+\cdots+T^{n-1}$, one has
\[
 [n]^\#(l)=l^n,
 \qquad [n]^\#(z)=N_n(l)z.
\]
\end{lemma}

\begin{proof}
On points, the first formula says that the $z$-coordinate of a product is
$z(g)+l(g)z(h)$, and the second has the opposite skewness.  Repeating the
same element $n$ times gives the geometric series.  Equivalently, the same
statement follows by induction from the coproduct formulas.
\end{proof}

When $l=1+t$ and $t^2=0$, the norm has the closed form
\begin{equation}\label{eq:binomial-norm}
 N_n(1+t)=n+\binom n2t.
\end{equation}
Apply this with $l=\lambda=1+\theta$.  Since $4=0$ and
$2\theta=2xV$,
\begin{equation}\label{eq:N4}
 N_4(\lambda)=4+6\theta=2\theta=2xV.
\end{equation}
Lemma~\ref{lem:norm} therefore gives
\begin{align}
 [4]^\#(U)&=2xVU=2xW,\nonumber\\
 [4]^\#(V)&=2xV^2=2xy^2V=4V=0,
 & [4]^\#(W)&=0.\label{eq:fourth}
\end{align}
The first element is nonzero: $2x\ne0$ by the Smith-normal-form calculation,
and $W$ is part of an $R$-basis of $A$.  Thus $[4]$ is not the unit word.

Also $\lambda^4=1$ and
\[
 N_8(\lambda)=N_4(\lambda)(1+\lambda^4)=2N_4(\lambda)=0.
\]
Since $A$ is generated by $U,V$, this says $[8]^\#=\iota\eps$.

For comparison, direct contraction of the coproduct gives the triangular
doubling calculation
\begin{equation}\label{eq:doubling}
 [2]^\#(U)=x^2yV-xW,
 \qquad [2]^\#(V)=-yW,
 \qquad [2]^\#(W)=0.
\end{equation}
Applying it twice leaves only the path
\[
 U\xmapsto{\ x^2y\ }V\xmapsto{\ -y\ }W,
 \qquad -(x^2y)y=-x^2y^2=2x.
\]
This recovers $[4]^\#(U)=2xW$ and makes visible the final socle carry.

\section{The ambient solvable group}

The two coproduct lines are standard coordinates on a small solvable group.
Let $\mathcal H$ have points $(l,u,v)$, with $l$ invertible, and law
\begin{equation}\label{eq:H-law}
 (l,u,v)(l',u',v')
  =(ll',\;u+lu',\;vl'+v').
\end{equation}
If $w=v/l$, this becomes
\[
 (l,u,w)(l',u',w')
  =(ll',\;u+lu',\;w+l^{-1}w').
\]
Thus
\[
 \mathcal H\simeq\Ga^2\rtimes\Gm,
\]
where $\Gm$ has weights $+1$ and $-1$.  It embeds in $\operatorname{GL}_3$
by
\[
 (l,u,w)\longmapsto
 \begin{pmatrix}
 l&0&u\\
 0&l^{-1}&w\\
 0&0&1
 \end{pmatrix}.
\]

The group $G$ is the finite closed subgroup of $\mathcal H$ cut out by the
two quadrics in \eqref{eq:A} and the graph equation
\begin{equation}\label{eq:graph}
 l=(1+yu)(1+xv)=1+yu+xv+xyuv.
\end{equation}
The group-like calculation in Section~4 says precisely that this graph and
the two quadrics are closed under \eqref{eq:H-law}.

In $\mathcal H$ one has the general power formula
\[
 (l,u,v)^n=(l^n,N_n(l)u,vN_n(l)).
\]
The counterexample is therefore the familiar semidirect-product norm
phenomenon.  Although $\lambda^4=1$, its translation norm is
\[
 N_4(\lambda)=2(\lambda-1)\ne0.
\]
The coefficient $\binom42=6\equiv2\pmod4$ is the characteristic-four
binomial carry.  The finite rank-four equations arrange that multiplying
this carry by $U$ gives exactly the top socle class $2xUV$.

This is a more precise description than calling the construction a
``Zappa--Szep-type product.''  No factorization into two finite flat subgroup
schemes is needed or asserted.

\section{A useful nonnormal rank-two subgroup}

There is nevertheless a useful rank-two subgroup
\[
 H_0=(V=0)\subset G,
 \qquad
 \mathcal O(H_0)=R[U]/(U^2-xyU).
\]
The ideal $(V)$ is a Hopf ideal: one has
$\Delta(V)=V\ot\lambda+1\ot V$, $\eps(V)=0$, and
$S(V)=V+yW=V(1+yU)\in(V)$.  Modulo $V$ one has
$\lambda=1+yU$, and a direct norm calculation gives $[2]_{H_0}=e$.

This subgroup is not normal.  On points, write $h=(a,0)\in H_0(T)$ and
$g=(u,v)\in G(T)$, suppressing the character coordinate.  The law
\eqref{eq:H-law} gives
\begin{equation}\label{eq:conjugation}
 V(ghg^{-1})=yva\,\lambda(g)^{-1},
\end{equation}
which is universally nonzero.  Indeed, its coordinate is
$(yV\lambda^{-1})\otimes a$, and
\[
 yV\lambda^{-1}=yV-y^2W\ne0
\]
in the free basis of $A$.  Thus the apparent rank-two-by-rank-two
filtration is not a filtration by normal subgroup schemes.

The fourth-power word factors nontrivially through $H_0$, because
$[4]^\#(V)=0$ but $[4]^\#(U)=2xUV\ne0$.  Squaring once more kills its image
because $H_0$ is killed by two.  This gives a structural explanation of
\[
 [4]\ne e,\qquad[8]=e,
\]
and pinpoints why a naive exponent argument based on two rank-two pieces
fails: the available rank-two subgroup is not normal.

\section{How minimal and how simple is the example?}

\subsection{Quotient-minimality}

Every nonzero ideal in an Artin local ring meets the socle.  Since
$\Soc(R)=\F_2(2x)$, every nonzero ideal of $R$ contains $2x$.  It follows
that every proper quotient of this particular base kills the coefficient
of $[4]^\#(U)$.  In particular,
\[
 R_8=R/(2x)
\]
has length eight, the kernel $\F_2(2x)$ is square-zero, and the fourth-power
defect vanishes over $R_8$.  The example is literally a one-dimensional
socle deformation that creates the defect.

\subsection{Minimality in the group-like bridge family}

There is also a precise, but restricted, universality statement.  Let $S$
be local, let $x,y\in\n_S$, and use over $S$ the same two quadrics, the same
$\lambda$, and the same skew-primitive formulas.  Put $q=xy^2$.  Coefficient
comparison in the two relation defects and in
$\Delta(\lambda)-\lambda\ot\lambda$ includes the necessary equations
\begin{equation}\label{eq:forcing}
 (q+1)(q+2)=0,
 \qquad y^3(q+1)=0,
 \qquad x^3(q+1)(q^2+q+1)=0.
\end{equation}
Because $q\in\n_S$, the factors $q+1$ and $q^2+q+1$ are units.  Hence the
Hopf formulas force
\[
 q=-2,\qquad x^3=y^3=0.
\]
Conversely, these relations imply $4=0$ and make the hand computations in
Section~4 go through.  Thus $R$ is universal for this displayed bridge
family.

If the fourth-power coefficient survives over such an $S$, the induced map
$R\to S$ is injective.  Indeed, a nonzero kernel would contain the socle
generator $2x$ and would kill $2xUV$.  Therefore a finite $S$ in this
family has at least $512$ elements; if its residue field is $\F_2$, it has
length at least nine.

This is not a global classification of all rank-four Hopf algebras.  The
correct scope is
\[
 \boxed{\text{length nine is minimal in this bridge family; global length
 eight remains open.}}
\]

\subsection{Generator-minimality}

Equation \eqref{eq:grR} gives
\[
 \dim_{\F_2}\m/\m^2=2.
\]
Since $2=xy^2\in\m^3$, the two base directions $x,y$ cannot be replaced by
one generator as a local $\Z/4$-algebra.  Likewise the special fiber of $A$
has two-dimensional cotangent space, so two generators are necessary as an
augmented coordinate algebra.
The displayed presentation is therefore already essentially
generator-minimal.  The genuine simplification is conceptual:
\[
 \lambda=(1+yU)(1+xV),
 \quad\text{two skew-primitive formulas, and one norm lemma.}
\]

\section{Independent verification and audit caveats}

The bundled archive \path{rank4_length9_counterexample_bundle.zip} contains
two bounded exact Python verifiers.  On the local Mac, both completed with
exit status zero under Python~3.14.5.  After extracting the archive, the
reproduction commands are
\begin{verbatim}
unzip -q rank4_length9_counterexample_bundle.zip \
  -d /tmp/rank4_length9_audit
cd /tmp/rank4_length9_audit
python3 verify_rank4_length9_counterexample.py
python3 independent_verify_rank4_length9_counterexample.py
\end{verbatim}
They verify the $512$-element base model, associativity, freeness,
multiplicativity and coassociativity of $\Delta$, the group-like identity,
the antipode equations, noncocommutativity, and the power words through
$[8]$.  The first completed in $3.60$ seconds and the second in $7.61$
seconds.  The archive passed \texttt{unzip -t}; its SHA-256 digest is
\begin{center}
\small\path{132f044bc092d1b888e5ff0e315e35a9b4790a92fd92809709d22448de0fe165}.
\end{center}

As a genuinely literal ideal check, from the workspace root we also ran
\begin{verbatim}
M2 --script m2/audit_rank4_length9_counterexample_20260711.m2
\end{verbatim}
Under Macaulay2~1.25.11 on the local Mac (not an SSH session), this bounded
computation completed with exit status zero in $0.23$ seconds.  Its
strong Groebner basis for the base ideal is
\[
 (4,\;2y,\;y^3,\;xy^2+2,\;2x^2,\;x^3),
\]
and it leaves $2$, $2x$, and $2xUV$ in nonzero normal form.  It independently
reduces both coproduct relation defects,
$\Delta(\lambda)-\lambda\ot\lambda$, the antipode defects, and the
$[8]$ coordinates to zero, while reducing $[4]^\#(U)$ to $2xUV$.

For the deliberately unoptimized construction, we separately ran
\begin{verbatim}
M2 --script m2/audit_rank4_unoptimized_bridge_20260711.m2
\end{verbatim}
over the untruncated five-coefficient ring
\eqref{eq:large-base}.  It completed with exit status zero.  A strong
Groebner reduction leaves $2s$ nonzero; it also verifies all five
compatibility relations and the nonvanishing of $2s$ in the
three-coefficient specialization \eqref{eq:three-coefficient-base}.
The same audit reduces both quadratic coproduct
defects and $\Delta(\lambda)-\lambda\ot\lambda$ to zero, verifies the
antipode identities, and returns
\[
 ([4]^\#(U),[4]^\#(V),[8]^\#(U),[8]^\#(V))
 =(2sUV,0,0,0).
\]

There are four minor caveats in the bundled write-up, none fatal.
\begin{enumerate}
\item The locality proof should explicitly use nilpotence of $x,y$ to prove
uniqueness of the maximal ideal; this was supplied in Section~2.
\item The two Python verifiers use different multiplication implementations,
but both begin from the same asserted seven-term normal form.  The Smith
normal form above and the independent literal ideal computation remove that
shared assumption.
\item The optional script
\path{derive_length9_family_conditions.py} requires SymPy, which was absent
from the default Python environment and was not listed as a dependency.  It
is not needed to certify the example.
\item ``Zappa--Szep-type'' should be read only as a heuristic.  The rigorous
structural statement is the closed embedding into
$\Ga^2\rtimes\Gm$ from Section~6.
\end{enumerate}

\section{Conclusion}

The example is valid, and its irreducible data fit in four lines:
\[
 \boxed{
 \begin{gathered}
 R=(\Z/4)[x,y]/(x^3,y^3,xy^2-2),\\
 A=R[U,V]/(U^2-xyU-x^2V,\;V^2-y^2V),\\
 \lambda=(1+yU)(1+xV),\\
 \Delta(U)=U\ot1+\lambda\ot U,
 \qquad \Delta(V)=V\ot\lambda+1\ot V.
 \end{gathered}}
\]
The conceptual mechanism is the semidirect-product norm
\[
 N_4(1+\theta)=4+\binom42\theta=2\theta.
\]
Here the bridge equations make $2\theta=2xV$, and multiplication by $U$
turns it into the nonzero top class $2xUV$.  This proves
$[4]\ne e$; the same norm vanishes at eight.  The result is smaller and much
more transparent than the earlier length-ten construction, but it does not
settle the global length-eight frontier.

\appendix
\section{Coefficient list for the universal bridge claim}

For completeness, here is the coefficient-level form of the
restricted universality calculation used in Section~8.  With $q=xy^2$, put
\[
 E=q^2+2q+2,
\]
\[
 F=x^3y^6+2x^2y^4+4q+2,
 \qquad
 H=x^3y^6+x^2y^4+q+2.
\]
After expanding preservation of both quadrics, group-likeness of $\lambda$,
and coassociativity in the free basis $1,U,V,UV$ of each tensor factor, the
primitive coefficients, up to sign and repetition, are
\begin{center}
\renewcommand{\arraystretch}{1.15}
\begin{tabular}{@{}ll@{}}
\toprule
$x^3y$ & $x^2y(q+2)$\\
$y^3(q+1)$ & $(q+1)(q+2)$\\
$xyE$ & $qE$\\
$x^2yE$ & $x^2y^2E$\\
$x^3y(q+2)^2$ & $y(q+1)^2(q+2)$\\
$x(q+1)^2(q+2)$ & $x^3y^2(q+1)^2$\\
$x^3y^3(q+1)^2$ & $x^2yF$\\
$x^3(q+1)(q^2+q+1)$ & $(q+1)H$\\
\bottomrule
\end{tabular}
\end{center}
The three entries isolated in \eqref{eq:forcing} force
$q=-2$, $y^3=0$, and $x^3=0$ over a local base.  Conversely, substitution of
these three relations annihilates every entry in the table.  This proves the
claimed ``if and only if'' inside this particular ansatz, without suggesting
that the ansatz classifies all rank-four group schemes.

\end{document}
